	\begin{algorithm} %[h]
		\SetKwFunction{algoDTestsFDSCALING}{} 
		\SetKwProg{myalg}{Algorithm}{}{} %write in 2nd agrument <<Algorithm>>, <<Procedure>> etc
		\nonl\myalg{\algoDTestsFDSCALING}{
			\KwInput{\\$S$ - количество частиц в рое, $n$ - количество итераций,\\
					 $\omega_0$ - начальное значение инерционного веса,\\
					 $\varphi_{p_{0}}$ - начальное значение когнитивного параметра,\\ $\varphi_{g_{0}}$ - начальное значение социального параметра,\\
					 $[t_{\min}, t_{\max}]$  - диапазон возможных значений фазы.
			}
			\KwOutput{ Оптимальное значение $f(g)$.}
			
			\For{каждой частицы $i = 1,\ldots, S$} {
				Инициализировать нач. значение позиции частицы равномерно распределённым случайным вектором $x_i \sim U(t_{\min}, t_{\max})$\;
				Инициализировать лучшую позицию из нач. приближения $p_i \leftarrow x_i$\;
				\If{$f(p_i) < f(g)$} {
					Обновить лучшую позицию среди роя $g \leftarrow p_i$\;
				}
				Инициализировать скорость частицы $v_i \sim U (- |t_{\max} - t_{\min}|, |t_{\max} - t_{\min}|)$\;
			}
			
			\While {не выполнено условие остановки} {
				\For{каждой частицы $i = 1,\ldots, S$}{
					\For{каждой размерности $d = 1, \ldots, k$} {
						Сгенерировать случайные числа $r_p, r_g \sim U(0, 1)$\;
						Обновить скорость частицы: $v_{id} \leftarrow \omega v_{id} +
						\varphi_p r_p (p_{id}-x_{id}) + \varphi_gr_g (g_d-x_{id})
						$\;
					}
					Обновить позицию частицы $x_i \leftarrow x_i + v_i$\;
				}
				\If{$f(x_i) < f(p_i)$}{
					Обновить лучшую позицию частицы $p_i \leftarrow x_i$\;
					\If{f($p_i$) < $f(g)$}{
						Обновить лучшую позицию среди роя $g \leftarrow p_i$\;
					}
				}
						
			}		
		}
		\caption{Псевдокод метода роя частиц \cite{Naidenova2017}}\label{alg:AlgoFDSCALING}
	\end{algorithm} 
	